\clearpage
\VerdierPhysicalPageAnchor{0154}%
\begingroup
\small
\setlength{\parskip}{0pt}%
\setlength{\emergencystretch}{1.5em}%

\VerdierSemanticAnchor{ch3-3-2}
\subsection{圏 \texorpdfstring{$\mathrm A(\mathcal D)$}{A(D)} の性質
  （\texorpdfstring{$\mathcal D$}{D} は三角圏）。}

\VerdierSemanticAnchor{ch3-3-2-0}
\noindent\textbf{3.2.0.}\quad
$\mathcal D$ と $\mathcal D'$ を二つの三角圏とする。
$\mathcal E x(\mathcal D,\mathcal D')$ を、対象が $\mathcal D$ から
$\mathcal D'$ への完全関手であり、射が完全関手の射である圏とする。同様に
$\mathbb Z\text{-}\operatorname{Add}(\mathrm A(\mathcal D),
\mathrm A(\mathcal D'))$ を、対象が $\mathrm A(\mathcal D)$ から
$\mathrm A(\mathcal D')$ への加法 $\mathbb Z$-関手であり、射が
$\mathbb Z$-関手の射である圏とする。以上により関手

\VerdierSemanticAnchor{ch3-3-2-0-1}
\[
\mathrm A_{\mathcal D,\mathcal D'}\colon
\mathcal E x(\mathcal D,\mathcal D')
\longrightarrow
\mathbb Z\text{-}\operatorname{Add}
(\mathrm A(\mathcal D),\mathrm A(\mathcal D')),
\tag{3.2.0.1}
\]
を定義した。これは完全関手 $F$ に $\mathrm A(F)$ を対応させる。

\VerdierSemanticAnchor{ch3-3-2-1}
\medskip
\noindent\textbf{定理 3.2.1.}\begingroup\renewcommand{\thefootnote}{(*)}%
\footnote{この定理は Freyd [17] から教示を受けたものである。}\endgroup\quad
\textit{a) $\mathcal D$ が三角圏なら、$\mathrm A(\mathcal D)$ はアーベル圏
である。$F\colon\mathcal D\to\mathcal D'$ が三角圏の間の完全関手なら、
$\mathrm A(F)$ (3.1.4.1) は完全である。関手
$\mathrm A_{\mathcal D,\mathcal D'}$ (3.2.0.1) は充満忠実かつ対象上単射である。}

\VerdierSemanticAnchor{ch3-3-2-1-b}
\textit{b) $\mathcal D$ の対象は $\mathrm A(\mathcal D)$ において入射対象かつ
射影対象である。$\mathrm A(\mathcal D)$ の任意の対象は、$\mathcal D$ の
ある対象の商対象（それぞれ部分対象）と同型である。}

\VerdierSemanticAnchor{ch3-3-2-1-c}
\textit{c) $\mathcal A$ をアーベル圏とする。
$\mathcal E x(\mathrm A(\mathcal D),\mathcal A)$ を、対象が
$\mathrm A(\mathcal D)$ から $\mathcal A$ への完全関手であり、射が関手の射である
圏とする。同様に $\mathcal H\!om(\mathcal D,\mathcal A)$ を、対象が
$\mathcal D$ から $\mathcal A$ への関手であり、射が関手の射である圏とする。
$\mathcal H\colon\mathcal D\to\mathrm A(\mathcal D)$ との合成は関手}
\[
\mathcal E x(\mathrm A(\mathcal D),\mathcal A)
\longrightarrow
\mathcal H\!om(\mathcal D,\mathcal A).
\]
\textit{を定める。この関手は充満忠実で、その本質像は $\mathcal D$ から
$\mathcal A$ へのコホモロジー的関手からなる (1.1.5)。特に
$\mathcal H\colon\mathcal D\to\mathrm A(\mathcal D)$ はコホモロジー的
関手である。}

\VerdierSemanticAnchor{ch3-3-2-1-proof-plan}
この定理の証明は段落 (3.2.9) まで続く。まず三つの中間命題を証明する。

\VerdierSemanticAnchor{ch3-3-2-2}
\medskip
\noindent\textbf{3.2.2.}\quad
$\mathcal D$ が $U$-小、すなわち $\mathcal D$ が $U$ の要素となるような
宇宙 $U$ をとる。$\mathcal D^\wedge$ を、
$\mathcal D$ 上の $U$-小アーベル群の圏に値をとる反変加法関手の圏とする。
$\mathcal D^\wedge$ はアーベル圏であり、$\mathcal D$ の各対象をそれが表現する関手へ
送る標準関手 $h\colon\mathcal D\to\mathcal D^\wedge$ は充満忠実である。
$h$ は $\mathcal D$ の各対象を $\mathcal D^\wedge$ の射影対象へ送る。
\endgroup
